\section{Returning to the neutral sector}\label{a:sec:neutral}

The deficit estimate was proved for ground states whose energy is within a
fixed distance of the minimum over all particle numbers.  We first apply it
at that minimum.  A screened deletion comparison will then put the neutral
and singly ionized ground states in the same class.

We record explicitly the harmonic-field estimate used in the deletion and
again in the final tail argument.  For a configuration $\omega=(x_i,\sigma_i)_{i=1}^n$
and $h>0$, write
\begin{equation}\label{a:eq:tail-inner-field}
 B_h(x;\omega)=V(x)-\sum_{|x_i|<h}K(x-x_i).
\end{equation}
This field is harmonic on $\{|x|>h\}$, configuration by configuration.

\Needspace{4\baselineskip}
\begin{lemma}[Conditional fields on annuli]\label{a:lem:tail-harmonic}
Fix positive constants $\alpha<\alpha'<\beta'<\beta$ and
$\alpha_0<\beta_0$.  Let $u>0$, $0<h\le\alpha u/2$, and let $\psi$ be a
normalized $\delta$-state with $n\le3Z$, where $\delta<\infty$ is arbitrary.
Let $Y$ be the data from a cut whose out support lies in
$\{\alpha_0u<|x|<\beta_0u\}$ and whose squared gradient cost is at most
$C u^{-2}$.  The cut may also be absent, in which case $Y$ is trivial.
There is a jointly measurable version
\[
 H_Y(x)=\E_\psi[B_h(x)\mid Y]
\]
which is harmonic on $\{\alpha u<|x|<\beta u\}$ for almost every datum.
Put
\[
 M(Y)=\sup_{\alpha'u\le|x|\le\beta'u}(H_Y(x))_+.
\]
Then
\begin{equation}\label{a:eq:tail-harmonic-bound}
 \|M\|_2\le C\left\{
 \frac{\max(u^{-3},1)+\sqrt{\delta u}}{u}
 +\sum_{j=0}^\infty
 \frac{\max((2^jh)^{-3},1)+\sqrt{\delta\,2^jh}}
      {\max(u,2^jh)}\right\}.
\end{equation}
The constant depends only on the fixed annular ratios and cut-gradient
constant, in addition to the fixed screening constants $L,A$.  In
particular it does not depend on $n,Z,\delta,h,u$.
\end{lemma}

\begin{proof}
For $x$ in the broad annulus, every source with $|x_i|<h$ lies outside
$B(x,Aa_x)$: indeed $h\le|x|/2$ and $A/L<1/2$.  Consequently
\begin{equation}\label{a:eq:tail-addback}
 B_h(x)=J_x+
 \sum_{\substack{|x_i|\ge h\\|x_i-x|\ge Aa_x}}K(x-x_i).
\end{equation}
If $2^jh\le|x_i|<2^{j+1}h$ and the summand occurs, then
\[
 K(x-x_i)\le\frac{C}{\max(u,2^jh)}.
\]
For radii comparable to or below $u$ this follows from the exclusion
$|x_i-x|\ge Aa_x\ge A\alpha u/L$; for radii greater than $2\beta u$ it
follows from $|x_i-x|\ge|x_i|/2$.

The initial-cut hypotheses of Proposition~\ref{a:prop:screen-local} hold
uniformly here.  Cover its support by a fixed finite number of cells
$B_z$ with $a_z\asymp u$, the constants being allowed to depend on $L$
and the annular ratios.  For each target $x$ in the broad annulus these
cell scales are comparable to $a_x$.  The gradient bound can be written
as $\sum_zD_z^2a_z^{-2}\1_{B_z}$, with $D_z^2\le C(a_z/u)^2$.
Since $a_zm_z\ge1$ and the screening parameter $P\ge1$,
\[
 \frac{D_z^2}{a_zm_z}\le C\le C_0\sqrt P.
\]
Enlarging the fixed cover constant $C_0$ suffices for every finite
$\delta$.
Thus \eqref{a:eq:screen-local}, conditional Jensen, the triangle inequality
in $L^2$, and \eqref{a:eq:screen-annular} applied to \eqref{a:eq:tail-addback}
give, pointwise in the spatial variable,
\[
 \|(H_Y(x))_+\|_2\le C\left\{
 \frac{\max(u^{-3},1)+\sqrt{\delta u}}{u}
 +\sum_{j\ge0}
 \frac{\max((2^jh)^{-3},1)+\sqrt{\delta\,2^jh}}
      {\max(u,2^jh)}\right\}.
\]
The series converges: its large-$j$ terms are bounded by a constant
times $2^{-j}h^{-1}+\sqrt\delta\,2^{-j/2}h^{-1/2}$.

For completeness, conditional fields here have simultaneous spatial
versions.  On every compact subset of $\{|x|>h\}$, the kernels in
\eqref{a:eq:tail-inner-field} and each fixed number of their derivatives
have deterministic bounds for the fixed finite system.  Conditional
integration therefore gives continuous harmonic functions, with
measurable dependence on the data.  One may first use the conditional
law of configurations supplied by the cut and then differentiate under
that integral.  For a harmonic function $H$, its positive part is
subharmonic.  Balls of radius $c u$ about points of the smaller annulus
remain in the broad annulus, so their mean-value inequalities imply
\[
 \sup_{\alpha'u\le|x|\le\beta'u}H(x)_+
 \le C u^{-3}\int_{\alpha u<|z|<\beta u}H(z)_+\,\dd z.
\]
Minkowski's integral inequality and the pointwise bound prove
\eqref{a:eq:tail-harmonic-bound}.  Suprema can be taken on countable dense
sets, which also proves their measurability.
\end{proof}

Two forms of this estimate will be used.  With fixed $t>0$, $h=t/4$,
$u=2^\ell t$, and $\delta=0$, it gives for the positive conditional
supremum on $u\le|x|\le2u$ the bound
\begin{equation}\label{a:eq:tail-deletion-field}
 \|M\|_2\le C_t\frac{1+\ell}{u}.
\end{equation}
Indeed the terms with $2^jh\le u$ have total numerator at most
$C_t(1+\ell)$, and the remaining series is $O_t(u^{-1})$.
With $h=u/8$, $u\ge1$, and the smaller annulus $u/2\le|x|\le4u$,
the broad annulus $u/4<|x|<6u$ is admissible and gives
\begin{equation}\label{a:eq:tail-distant-field}
 \|M\|_2\le C\big(u^{-1}+\sqrt\delta\,u^{-1/2}\big).
\end{equation}
In the first application the broad annulus can be $u/2<|x|<3u$.
The non-strict equality $h=\alpha u/2$ causes no boundary issue, since
the broad annulus is open and still separated from the sources.

\begin{lemma}[An upper bound on an inner deficit]\label{a:lem:tail-deficit-upper}
For each fixed $t>0$ and every normalized $\delta$-state with $n\le3Z$,
\begin{equation}\label{a:eq:tail-deficit-upper}
 Z-\E\#\{i:|x_i|<t/4\}\le C_t(1+\sqrt\delta).
\end{equation}
\end{lemma}

\begin{proof}
Use Lemma~\ref{a:lem:tail-harmonic} with no cut, $h=t/4$, and $u=t$.
Its convergent series is bounded by $C_t(1+\sqrt\delta)$.  The
spherical average of the expected field $B_{t/4}$ on $|x|=t$ equals
\[
 \frac{Z-\E\#\{i:|x_i|<t/4\}}{t},
\]
by the spherical mean of the Coulomb kernel.  Bounding this average
from above by the positive supremum proves the assertion, with the
fixed factor $t$ absorbed into $C_t$.
\end{proof}

\begin{proposition}[A bounded offset for neutral sectors]\label{a:prop:tail-offset}
There are universal constants $C,Z_0<\infty$ such that, for $Z\ge Z_0$,
\begin{equation}\label{a:eq:tail-offset}
 E_{Z-1}\le E+C.
\end{equation}
\end{proposition}

\begin{proof}
Take a normalized ground state $\Psi_*$ in sector $N_*$, where
$Z\le N_*\le3Z$ and $E_{N_*}=E$, as provided by
Lemma~\ref{lem:atomic-sector}.  Proposition~\ref{a:prop:prop-deficit}, with
$D_0=0$, supplies a fixed sufficiently small $t>0$ such that
\[
 \E_{\Psi_*}\big[Z-\#\{i:|x_i|<2t\}\big]>1
\]
for all sufficiently large $Z$.  Choose a smooth square partition
$c^2+o^2=1$, with $c=1$ on $|x|\le t$, $c=0$ on $|x|\ge2t$, and
$|\nabla c|^2+|\nabla o|^2\le C t^{-2}$ supported in the intervening
annulus.  Let $X$ be its full out data and set
$q^c=Z-n_c$, where $n_c$ is the number of core particles.  Since
$n_c\le\#\{i:|x_i|<2t\}$, we have $\E q^c>1$.

We first show that the favorable event $G=\{n_c\le Z-1\}$ has
probability bounded below independently of $Z$.  Independently of the
cut labels, observe one unordered array with the smooth compact noise
law of Lemma~\ref{lem:obs-tilt}, choosing its fixed width so that each
spatial displacement is less than $t/8$.  Put
\[
 q^{\rm obs}=Z-\#\{\text{array entries in }B(0,t/2)\}.
\]
Every entry counted here comes from a particle of radius less than
$5t/8$, which must be core.  Thus $q^c\le q^{\rm obs}$ on the enlarged
probability space.  Conversely every particle of radius less than
$t/4$ produces an entry in $B(0,t/2)$.

For an event $A_\lambda=\{q^{\rm obs}>\lambda\}$ of probability
$p>0$, the single-array version of \eqref{eq:obs-tilt} produces a
normalized state $\Psi_{A_\lambda}$ whose raw configuration law is the
law conditioned on that event and whose offset satisfies
\[
 \delta\le C_t\log^5(e/p).
\]
Every compatible configuration obeys
$Z-\#\{|x_i|<t/4\}>\lambda$.  Hence
Lemma~\ref{a:lem:tail-deficit-upper} gives
\[
 \lambda\le C_t\big(1+\log^{5/2}(e/p)\big).
\]
For sufficiently large $\lambda$ this implies
$p\le e\exp(-c_t\lambda^{2/5})$; the inequality is also harmless when
$p=0$.  Integrating the tail yields
\begin{equation}\label{a:eq:tail-favorable-probability}
 \E(q^c)_+^2\le\E(q^{\rm obs})_+^2\le C_t,
 \qquad
 \Prob(G)=\Prob(q^c>0)
 \ge\frac{(\E(q^c)_+)^2}{\E(q^c)_+^2}\ge c_t>0.
\end{equation}
Here the equality of events uses the integer-valued core count, and
$\E(q^c)_+\ge\E q^c>1$.

We next estimate the average energy of the core slices.  The
localization error is $O_t(1)$ by \eqref{a:eq:screen-annular} at offset
zero.  In the attractive term of \eqref{eq:loc-cut}, retain only
the core repulsion from particles with radius less than $t/4$, which
are core with probability one.  Conditional integration followed by
averaging the original cut labels gives
\begin{equation}\label{a:eq:tail-deletion-attraction}
 \E\sum_{i\ {\rm out}}\Phi_X(x_i)
 \le\E_{\Psi_*}\sum_i o(x_i)^2 B_{t/4}(x_i).
\end{equation}
This comparison of expectations uses the raw inner particles; it does
not identify fields from different core slices.

Divide the nonzero terms on the right into shells
$u\le|x_i|<2u$, $u=2^\ell t$, $\ell\ge0$.  For each shell use an
auxiliary cut of the original state, full out on that shell, supported
in $u/2<|x|<4u$, with gradient cost $C u^{-2}$, and call its data $Y$.
All weighted shell positions and their weights $o(x_i)^2$ are then
recorded.  Conditional integration of the raw field, with the harmonic
version of Lemma~\ref{a:lem:tail-harmonic}, bounds the shell expectation
by $\E[n_u M(Y)]$, where $n_u=\#\{u\le|x_i|<2u\}$.
Equations~\eqref{a:eq:screen-annular} and
\eqref{a:eq:tail-deletion-field} give
\[
 \E[n_uM]\le\|n_u\|_2\|M\|_2
 \le C_t\frac{1+\ell}{2^\ell t}.
\]
The constants are independent of the shell index and this series is
summable.  For each fixed system the original attraction and
repulsion expectations are finite by the form bounds, so the shell
decomposition is legitimate; alternatively one can first use finitely
many shells and pass to the limit.  Dropping the nonnegative out
kinetic energy and out--out repulsion from
\eqref{eq:loc-cut}, and using
\eqref{a:eq:tail-deletion-attraction}, now proves
\begin{equation}\label{a:eq:tail-deletion-energy}
 0\le\E(e_X-E)\le C_t.
\end{equation}

Every slice satisfies $e_X\ge E$, and on $G$ sector monotonicity gives
$e_X\ge E_{n_c}\ge E_{Z-1}$.  Therefore
\[
 c_t(E_{Z-1}-E)
 \le\Prob(G)(E_{Z-1}-E)
 \le\E(e_X-E)\le C_t,
\]
which proves \eqref{a:eq:tail-offset}.  This argument has used only the
zero-offset ground state and the general-offset screening estimates;
no neutral offset or ionization gap was assumed.
\end{proof}

By monotonicity, $E_Z\le E_{Z-1}\le E+C$.  We may consequently apply
Proposition~\ref{a:prop:prop-deficit} again, now with this fixed value of
$D_0$, to every normalized neutral ground state.  There are universal
$s_*>0$ and $Z_1$ such that
\begin{equation}\label{a:eq:tail-lower}
 \int_{|x|\ge s_*}\rho_{\Psi_Z}(x)\,\dd x>\frac12
 \qquad (Z\ge Z_1).
\end{equation}
In fact the deficit proposition gives a tail greater than one at an
appropriate such radius.  Absolute continuity of the density makes
the choice between strict and non-strict radial endpoints immaterial.
